% Options for packages loaded elsewhere
\PassOptionsToPackage{unicode}{hyperref}
\PassOptionsToPackage{hyphens}{url}
\PassOptionsToPackage{dvipsnames,svgnames,x11names}{xcolor}
\documentclass[
]{article}
\usepackage{xcolor}
\usepackage[margin=25mm]{geometry}
\usepackage{amsmath,amssymb}
\setcounter{secnumdepth}{-\maxdimen} % remove section numbering
\usepackage{iftex}
\ifPDFTeX
  \usepackage[T1]{fontenc}
  \usepackage[utf8]{inputenc}
  \usepackage{textcomp} % provide euro and other symbols
\else % if luatex or xetex
  \usepackage{unicode-math} % this also loads fontspec
  \defaultfontfeatures{Scale=MatchLowercase}
  \defaultfontfeatures[\rmfamily]{Ligatures=TeX,Scale=1}
\fi
\usepackage{lmodern}
\ifPDFTeX\else
  % xetex/luatex font selection
\fi
% Use upquote if available, for straight quotes in verbatim environments
\IfFileExists{upquote.sty}{\usepackage{upquote}}{}
\IfFileExists{microtype.sty}{% use microtype if available
  \usepackage[]{microtype}
  \UseMicrotypeSet[protrusion]{basicmath} % disable protrusion for tt fonts
}{}
\makeatletter
\@ifundefined{KOMAClassName}{% if non-KOMA class
  \IfFileExists{parskip.sty}{%
    \usepackage{parskip}
  }{% else
    \setlength{\parindent}{0pt}
    \setlength{\parskip}{6pt plus 2pt minus 1pt}}
}{% if KOMA class
  \KOMAoptions{parskip=half}}
\makeatother
\setlength{\emergencystretch}{3em} % prevent overfull lines
\providecommand{\tightlist}{%
  \setlength{\itemsep}{0pt}\setlength{\parskip}{0pt}}
\usepackage{bookmark}
\IfFileExists{xurl.sty}{\usepackage{xurl}}{} % add URL line breaks if available
\urlstyle{same}
\hypersetup{
  colorlinks=true,
  linkcolor={Maroon},
  filecolor={Maroon},
  citecolor={Blue},
  urlcolor={Blue},
  pdfcreator={LaTeX via pandoc}}

\author{}
\date{}

\begin{document}

\section{Resultant forms and
elimination}\label{resultant-forms-and-elimination}

\emph{Written by GPT-6.1 Sol (OpenAI) in Codex at Ultra reasoning
effort, October 2026. Self-checked by the writing AI, GPT-6.1 Sol, at
Ultra. No independent review is claimed. Public domain (CC0).}

A determinant can remember every point of a finite polynomial scheme,
together with its length. A second polynomial, the minimal polynomial,
measures how those points and nilpotents can be distinguished by one
generic linear function. Comparing the two leads to Noether's tests for
prime and primary ideals. We first build the finite-dimensional
construction, then retain all dimensional layers when passing to generic
fibres.

Let \(k\) be a field. Rings are commutative with identity. A
zero-dimensional ideal is proper, with a nonzero quotient of Krull
dimension zero. The proved prerequisites are
\href{../prerequisites/AG-CA/AG-CA-03.html}{Noetherian and Artinian
rings}, Theorem 4.2 for Artinian structure and product decomposition;
\href{../prerequisites/AG-CA/AG-CA-09.html}{Krull dimension and Noether
normalization}, Theorem 4.2 for dimension and transcendence degree;
\href{../prerequisites/AG-CA/AG-CA-04.html}{Associated primes and
primary decomposition}, Theorems 1.2, 3.1, 4.2 and 5.3; and the
preceding \href{../NOE-IDL-02.html}{primary-ideal lesson} for
contraction and primary localization. We use
\href{../prerequisites/AG-CA/AG-CA-06.html}{the Nullstellensatz},
Theorems 1.3 and 2.1--2.2, and
\href{../prerequisites/AG-CA/AG-CA-02.html}{localization}, Theorem 4.2
for Nakayama. Linear algebra and polynomial factorization are elementary
background. The finite-field and differential bridges needed here are
proved in Section 5. Our free historical edition contains Hentzelt's
work edited by Noether and Noether's \emph{Elimination Theory and
General Ideal Theory}, linked below.

\subsection{1. Fixing the sign of the
resultant}\label{fixing-the-sign-of-the-resultant}

For \(f=a_m z^m+\cdots+a_0\) and \(g=b_nz^n+\cdots+b_0\), with
\(m,n\ge1\), use descending-power bases in the map

\[
S:k[z]_{<n}\oplus k[z]_{<m}\longrightarrow k[z]_{<m+n},
\qquad (u,v)\longmapsto uf+vg.
\]

The first block is the \(f\)-block. Define
\(\operatorname{Res}(f,g)=\det S\). This specifies a column version of
Sylvester's matrix and its sign once and for all.

\textbf{Theorem 1.1.} If \(a=a_m\ne0\), \(b=b_n\ne0\), and in an
algebraic closure \(f=a\prod_{i=1}^m(z-\alpha_i)\),
\(g=b\prod_{j=1}^n(z-\beta_j)\), then

\[
\operatorname{Res}(f,g)=a^n b^m\prod_{i,j}(\alpha_i-\beta_j).
\tag{1}
\]

In particular it vanishes exactly when the polynomials have a common
root. Over the universal coefficient ring
\(C=\mathbb Z[a_0,\ldots,a_m,b_0,\ldots,b_n]\), there are
\(u\in C[z]_{<n}\), \(v\in C[z]_{<m}\) with

\[
uf+vg=\operatorname{Res}(f,g).
\tag{2}
\]

\textbf{Proof.} First take both polynomials monic with independent
roots. A common root makes evaluation annihilate the image of \(S\), so
the determinant vanishes when any \(\alpha_i=\beta_j\). In the
polynomial ring over \(\mathbb Z\) in all roots, each prime element
\(\alpha_i-\beta_j\) therefore divides the determinant; their
distinctness implies that their product divides it. Its degree in a
fixed \(\alpha_i\) is at most \(n\), since there are \(n\) columns from
\(f\), each at most linear in that root. Its degree in a fixed
\(\beta_j\) is at most \(m\). Thus the quotient by the product is a
constant integer.

To find the constant, set \(f=z^m\), \(g=(z-1)^n\). The \(f\)-columns
give the leading \(n\) coordinates as an identity block. The remaining
block is multiplication by \(g\) modulo \(z^m\), in descending powers,
triangular with diagonal \(g(0)=(-1)^n\). The determinant is
\((-1)^{mn}\), exactly the value of the product in (1); the constant is
one. Scaling the first and second blocks by \(a,b\) gives their powers
\(a^n,b^m\). This polynomial identity specializes to every field,
including repeated roots and positive characteristic.

Finally, \(S\operatorname{adj}(S)=(\det S)\mathrm{Id}\) over \(C\).
Apply it to the coordinate vector of the constant polynomial one. The
two blocks of \(\operatorname{adj}(S)1\) give \(u,v\), proving (2)
without dividing by either leading coefficient. \(\square\)

Swapping the two blocks proves
\(\operatorname{Res}(g,f)=(-1)^{mn}\operatorname{Res}(f,g)\). If one
polynomial is a nonzero constant, the corresponding convention is
\(\operatorname{Res}(a,g)=a^n\) and \(\operatorname{Res}(f,b)=b^m\).
Formula (2) is a universal identity for positive formal degrees; it does
not assert that the empty determinant one for two formal degree-zero
polynomials belongs to their universal ideal.

For example,

\[
\operatorname{Res}(z^2+bz+c,z-a)=a^2+ba+c.
\]

The discriminant of a degree-\(m\) polynomial with leading coefficient
\(a\) is \(a^{2m-2}\prod_{i<j}(\alpha_i-\alpha_j)^2\). Taking \(g=f'\)
in (1) gives

\[
\operatorname{disc}(f)=(-1)^{m(m-1)/2}a^{-1}\operatorname{Res}(f,f'),
\tag{3}
\]

where the derivative is placed in its formal degree-\(m-1\) block. This
is a polynomial identity after cancellation of \(a\), and remains valid
when the derivative's actual degree drops in positive characteristic.
For \(f=Az^2+Bz+C\), its resultant with \(2Az+B\) is \(-A(B^2-4AC)\),
giving the familiar discriminant. When a specialized leading coefficient
vanishes, the formal Sylvester determinant can also detect behaviour at
infinity; the common-root equivalence in Theorem 1.1 requires the stated
actual degrees.

\subsection{2. Multiplication sees values and
lengths}\label{multiplication-sees-values-and-lengths}

Let \(I\subset k[x_1,\ldots,x_n]\) be zero-dimensional and \(A=k[x]/I\).
It is finite-dimensional over \(k\): it is Noetherian Artinian, its
finitely many residue fields are finite over \(k\) by Zariski's lemma,
and a composition series has these fields as factors. Multiplication by
\(f\) is a \(k\)-linear operator \(M_f\).

Suppose first that \(k\) is algebraically closed. The Artinian product
decomposition gives

\[
A=\prod_{\xi\in V(I)}A_\xi,
\qquad m_\xi=\dim_k A_\xi.
\]

Here \(A_\xi\) is the localization at the point's maximal ideal. Its
maximal ideal is nilpotent and its residue field is \(k\); \(m_\xi\) is
its length.

\textbf{Theorem 2.1 (Stickelberger).}

\[
\det(T\mathrm{Id}-M_f)=\prod_{\xi\in V(I)}(T-f(\xi))^{m_\xi}.
\tag{4}
\]

\textbf{Proof.} On \(A_\xi\), the class of \(f\) is \(f(\xi)+h\) with
\(h\) in the nilpotent maximal ideal. Filter this factor by powers of
that ideal. On every successive vector-space quotient multiplication by
\(h\) is zero, so \(M_f\) acts as the scalar \(f(\xi)\). A basis
respecting the filtration gives a block triangular matrix with that
scalar throughout its diagonal. The determinant is
\((T-f(\xi))^{m_\xi}\). Multiplication respects the product factors, so
the full determinant is their product. \(\square\)

The theorem concerns algebraic multiplicity of eigenvalues, not the
dimension of their eigenspaces. On \(k[x,y]/(x^2,y)\), multiplication by
\(x\) is a nonzero nilpotent operator with characteristic polynomial
\(T^2\), while its zero eigenspace has dimension one.

For arbitrary \(k\), base change to \(\bar k\). The matrix entries and
determinant commute with base change. In (4), use the geometric points
and the dimensions of the local factors of \(A\otimes_k\bar k\). These
geometric dimensions incorporate inseparability as well as nilpotent
thickness.

\subsection{3. The resultant form}\label{the-resultant-form}

With independent variables \(u_0,\ldots,u_n\), define

\[
R_I(u)=\det\left(u_0\mathrm{Id}+\sum_{j=1}^nu_jM_{x_j}\right).
\tag{5}
\]

It is homogeneous of degree \(\dim_k A\), monic in \(u_0\), independent
of the chosen basis, and has coefficients in \(k\).

\textbf{Theorem 3.1.} Over \(\bar k\),

\[
R_I(u)=\prod_{\xi\in V_{\bar k}(I)}
\left(u_0+\sum_j u_j\xi_j\right)^{m_\xi}.
\tag{6}
\]

\textbf{Proof.} Apply the maximal-ideal filtration argument from Theorem
2.1 to multiplication by the generic linear function \(\sum u_jx_j\),
over \(\bar k(u_1,\ldots,u_n)\), and take \(T=-u_0\), with the
determinant signs cancelling. The resulting identity of polynomials is
(6). \(\square\)

Unique factorization and the normalization of each factor's
\(u_0\)-coefficient recover the point coordinates and their
multiplicities. They do not recover the ideal. For instance \((x^2,y)\)
and \((x,y^2)\) both have form \(u_0^2\), despite different nilpotent
directions.

\textbf{Theorem 3.2 (primary criterion).} A zero-dimensional ideal is
primary exactly when its geometric points form one orbit under
\(\operatorname{Aut}(\bar k/k)\). In that case every linear factor in
(6) has the same exponent. If \(\mathfrak m=\sqrt I\),
\(L=k[x]/\mathfrak m\), and \(\ell\) is the length of \(A\) as a local
ring, then

\[
R_I=R_{\mathfrak m}^{\ell}
=N_{L/k}\left(u_0+\sum_j u_j\bar x_j\right)^{\ell}.
\tag{7}
\]

Over an algebraically closed field the criterion says exactly that
\(R_I\) is a power of one linear form.

\textbf{Proof.} The quotient's Artinian factors are indexed by its
maximal ideals. A finite Artinian ring is primary as a zero ideal
exactly when it is local: in a local factor all nonunits are nilpotent;
two nonzero product factors supply zero divisors which are not
nilpotent. The embeddings of a finite residue field \(L\) into
\(\bar k\) form one orbit, and two different maximal ideals cannot lie
in the same orbit since their relation ideals over \(k\) are different.
Thus one orbit is equivalent to one maximal ideal.

For a local factor, a composition series has \(\ell\) copies of \(L\).
Multiplication preserves it and induces the same multiplication operator
on each copy, proving (7) by a triangular determinant. Write
\([L:k]=sd\), where \(s\) is the number of embeddings and \(d\) its
inseparable degree. After base change, each embedding contributes a
local algebra of dimension \(d\); equivalently the norm is the product
of its \(s\) linear forms, each to power \(d\). The common exponent in
(6) is therefore \(d\ell\). \(\square\)

The squarefree orbit product need not itself have coefficients in \(k\).
For \(k=\mathbb F_p(t)\), \(I=(x^p-t,y)\) is maximal and has local
length one, but \(R_I=(u_0+u_1t^{1/p})^p=u_0^p+t u_1^p\). Its sole
geometric point has multiplicity \(p\). This distinction corrects an
unqualified use of ``Galois orbit product'' over imperfect fields.

\subsection{4. A finite scheme computed by
matrices}\label{a-finite-scheme-computed-by-matrices}

Take \(I=(x^2-y,y^2-y)\). Eliminating \(y\) identifies its quotient with
\(k[x]/(x^4-x^2)\), of dimension four in every characteristic. With
basis \(1,x,x^2,x^3\), columns recording images of basis vectors, we
obtain

\[
M_x=\begin{pmatrix}0&0&0&0\\1&0&0&0\\0&1&0&1\\0&0&1&0\end{pmatrix},
\qquad
M_y=M_x^2=\begin{pmatrix}0&0&0&0\\0&0&0&0\\1&0&1&0\\0&1&0&1\end{pmatrix}.
\]

Their characteristic polynomials are \(T^2(T^2-1)\) and \(T^2(T-1)^2\).
If \(\operatorname{char}k\ne2\), the geometric points are
\((0,0),(1,1),(-1,1)\), with lengths \(2,1,1\). Consequently

\[
R_I=u_0^2(u_0+u_1+u_2)(u_0-u_1+u_2).
\tag{8}
\]

In characteristic two the last two points coincide:
\(x^4-x^2=x^2(x-1)^2\). The two points have length two each, and (8)
becomes \(u_0^2(u_0+u_1+u_2)^2\). Both the multiplication calculation
and its characteristic qualification matter.

\subsection{5. Retaining every dimension in a generic
fibre}\label{retaining-every-dimension-in-a-generic-fibre}

Localizing only at parameters of the largest component can discard
embedded components. We instead make the isolated-component filtration
from the second lesson explicit. Let \(P=k[x_1,\ldots,x_n]\),
\(I\subsetneq P\). Adjoin an invertible generic matrix \(U\), put
\(F=k(U)\), and set \(z=Ux\). Write \(I_F\) for the transformed ideal in
\(F[z]\).

We first record the finite-field calculations used for generic
parameters. They also justify the field and orbit assertions in Section
3.

\textbf{Finite-field calculation.} An irreducible polynomial over a
characteristic-zero field is separable: its derivative is nonzero of
smaller degree, so it is relatively prime to the polynomial. In
characteristic \(p>0\), every monic irreducible polynomial has the form
\(g(T^{p^e})\), with \(g\) irreducible and \(g'\ne0\). Repeatedly
extract powers of \(T^p\) until the derivative is nonzero. Its distinct
roots are the \(p^e\)-th roots of the distinct roots of \(g\), and their
multiplicities are \(p^e\).

An embedding of a field extends through a finite simple algebraic
extension by choosing a root of the transported minimal polynomial. For
a finite tower, the number of embeddings into an algebraic closure is
therefore the product of the numbers of distinct roots at the successive
steps. Each is the degree of the step divided by a power of \(p\). An
extension generated by separable elements has as many embeddings as its
degree, because at each step the minimal polynomial divides a separable
polynomial. Every element of that extension is separable: restriction to
its simple subfield, followed by extension through the remaining tower,
bounds the number of embeddings by the number of embeddings of that
subfield times the remaining degree. Equality with the full degree
forces the simple subfield to have as many distinct embeddings as its
degree. In particular, sums, products and inverses of separable elements
are separable.

In characteristic \(p>0\), if a finite extension has only one embedding,
each simple subfield has only one embedding, since all of its embeddings
extend. The preceding polynomial calculation then makes each element
purely inseparable, and the extension has degree \(p^g\). For a general
finite extension, the ratio of its degree to its number of embeddings is
also a power of \(p\), by the tower calculation. In characteristic zero
that ratio is one. Any two embeddings of a finite extension of \(k\)
into \(\bar k\) are conjugate by an automorphism of \(\bar k/k\): extend
the induced isomorphism between their images to an embedding of
\(\bar k\) in itself by a maximal partial extension. A missing algebraic
element could always be added by choosing a root, so the maximal domain
is all of \(\bar k\). Its image is algebraically closed and \(\bar k\)
is algebraic over it, so the embedding is onto. This proves the orbit
assertion used in Theorem 3.2. The resulting automorphisms of
\(A\otimes_k\bar k\) carry its local factors to one another and preserve
their dimensions. Applied to a residue field of degree \(sd\), with
\(s\) embeddings, this also gives dimension \(d\) to each of its \(s\)
geometric factors, as used in (7).

\textbf{Differential calculation.} We use only the construction,
polynomial and localization rules, and the first fundamental sequence
from \href{../prerequisites/AG-CA/AG-CA-16.html}{Kähler differentials},
Theorem 1.1, Theorems 2.2, 3.1 and 5.1. In particular,

\[
E\otimes_C\Omega_{C/K}\longrightarrow\Omega_{E/K}
\longrightarrow\Omega_{E/C}\longrightarrow0
\]

is exact for a tower \(K\subset C\subset E\). We now prove the
additional facts we need.

For a finite extension \(E/K\) of characteristic \(p\), set \(C=K E^p\),
the subfield generated by \(K\) and all \(p\)-th powers in \(E\). Every
\(K\)-derivation kills \(C\), so \(\Omega_{E/K}=\Omega_{E/C}\). Choose
generators \(b_1,\ldots,b_r\) over \(C\) successively, omitting a
generator already in the preceding field. Since \(b_j^p\in C\), each
remaining step has degree \(p\). Indeed \(T^p-b_j^p\) has only one
distinct root; its irreducible factor is either linear or has degree
\(p\), by the preceding polynomial calculation. The monomials
\(b_1^{e_1}\cdots b_r^{e_r}\), \(0\le e_j<p\), form a basis over \(C\).
Consequently

\[
E=C[T_1,\ldots,T_r]/(T_j^p-b_j^p:1\le j\le r),
\qquad \Omega_{E/C}=\bigoplus_{j=1}^r E\,db_j.
\]

The presentation is an isomorphism because its spanning monomials map to
that basis. Its relations have zero differential, proving the displayed
differential formula. Thus \(\Omega_{E/K}=0\) exactly when \(E=K E^p\).

This vanishing implies separability. Iterating \(E=K E^p\) gives
\(E=K E^{p^N}\) for every \(N\). If \(\alpha_1,\ldots,\alpha_m\)
generate \(E/K\), the irreducible-polynomial calculation makes a
sufficiently high \(p\)-power of each \(\alpha_j\) separable over \(K\).
For one common \(N\), therefore,
\(K E^{p^N}=K(\alpha_1^{p^N},\ldots,\alpha_m^{p^N})\) is separable, by
the finite-field calculation. Since it is \(E\), the original extension
is separable. Conversely a separable extension has zero differentials by
\href{../prerequisites/AG-CA/AG-CA-16.html}{Kähler differentials},
Theorem 6.2, whose derivation-extension proof is already present. We
have proved the finite criterion

\[
E/K\text{ finite is separable}\quad\Longleftrightarrow\quad
\Omega_{E/K}=0.
\tag{5a}
\]

Finally let \(L/k\) be finitely generated of transcendence degree \(d\),
with \(k\) perfect of characteristic \(p\). Choose any transcendence
basis \(t\), and put \(D=[L:k(t)]\). Frobenius identifies \(L\) with
\(L^p\) and \(k(t)\) with \(k(t)^p\), so \([L^p:k(t)^p]=D\). The
monomials in the \(t_j\) with exponents less than \(p\) form a basis of
\(k(t)\) over \(k(t)^p=k(t_1^p,\ldots,t_d^p)\). Independence follows by
clearing denominators and comparing distinct exponent classes modulo
\(p\); their span is a finite-dimensional domain, hence a field, and
contains all the \(t_j\). Computing \([L:k(t)^p]\) in two ways now gives

\[
D p^d=[L:L^p]D,\qquad [L:L^p]=p^d.
\]

Because \(k\subset L^p\), the same generator and presentation argument
over \(L^p\) supplies precisely \(d\) differential basis elements. Thus
\(\dim_L\Omega_{L/k}=d\). In characteristic zero this dimension follows
from Theorem 6.2 of the differential lesson: any transcendence basis
leaves a finite algebraic extension, which is separable by the first
paragraph. These arguments prove the dimension assertion without
assuming a separating basis exists.

We can now establish linear normalization and its separability
qualification.

\textbf{Lemma 5.1.} For each of the finitely many associated primes
\(\mathfrak p\) of \(P/I\), of dimension \(d\), the last \(d\) generic
coordinates are algebraically independent modulo \(\mathfrak p_F\), and
its function field is finite over \(F(z_{n-d+1},\ldots,z_n)\). Any list
of more than \(d\) coordinates satisfies a nonzero polynomial relation
there. If \(k\) is perfect, the finite extension for the generic list of
length \(d\) is separable.

\textbf{Proof of normalization.} Fix one prime, let \(A=P/\mathfrak p\),
and let \(L\) be its function field. We prove the generic assertion
directly, including finite coefficient fields. If \(n>d\), choose a
nonzero relation \(f(x)=0\), with highest homogeneous part \(f_D\).
Introduce independent coefficients \(c_2,\ldots,c_n\) and put
\(y_j=x_j-c_jx_1\). After substituting \(x_j=y_j+c_jx_1\), the
coefficient of \(x_1^D\) is the nonzero polynomial
\(f_D(1,c_2,\ldots,c_n)\). It is nonzero because homogeneity makes its
monomials have distinct exponent lists in the \(c_j\). Dividing by this
coefficient makes \(x_1\) integral over the algebra generated by
\(y_2,\ldots,y_n\). All the original generators are then integral, so
the inclusion is finite. Its smaller algebra is a domain of dimension
\(d\), by equality of dimension under finite integral inclusions,
\href{../prerequisites/AG-CA/AG-CA-05.html}{Integral extensions},
Theorem 5.1.

Repeat with fresh independent coefficients whenever the number of
remaining generators exceeds \(d\). After \(n-d\) steps the remaining
\(d\) generators are independent, by the dimension--transcendence-degree
theorem, and the full algebra is finite over their polynomial algebra.
All these changes are linear and triangular. The final retained
coordinates have the shape

\[
w_j=x_{n-d+j}+\sum_{h=1}^{n-d}a_{jh}x_h\quad(1\le j\le d).
\]

The \(d(n-d)\) coefficients \(a_{jh}\) are independent variables,
together with the coefficients used among the first \(n-d\) eliminated
coordinates: at the step eliminating the \(h\)-th coordinate, the new
coefficient in each retained coordinate is its fresh coefficient with
sign minus, plus an expression in later-step coefficients. Starting at
the last step and solving backwards recovers these fresh coefficients
from the \(a_{jh}\) and the coefficients among eliminated coordinates.
This is an invertible triangular change of the independent coefficient
variables. If the latter extra coefficients are denoted \(b\), we have
proved finiteness of the function field over \(k(a,b)(w)\). Adjoining
the independent \(b\) cannot reduce transcendence degree of
\(L(a)/k(a)(w)\); thus this finitely generated extension already has
transcendence degree zero, and is finite. When \(n=d\), the original
generators are independent and there is no elimination step.

For the last \(d\) rows of a completely generic matrix \(U\), their last
\(d\) columns form an invertible matrix \(B\). Multiplication by
\(B^{-1}\) writes these rows uniquely as \([a\mid\mathrm{Id}_d]\); their
first-block coefficients \(a\), the entries of \(B\), and the remaining
rows give independent rational coordinates on the matrix parameter
field. The assertion just proved for \(w\) consequently proves it for
the last \(d\) coordinates of \(U\). For \(d=0\), the argument says
simply that the residue field is finite over the coefficient field. The
argument works for each of the finite list of associated primes with the
same independent matrix \(U\); no choice of a rational matrix over \(k\)
is needed. More than \(d\) coordinates must be dependent by
transcendence degree.

\textbf{Proof of separability.} Assume \(k\) perfect, put \(M=L(U)\),
\(F=k(U)\), and write \(t\) for the last \(d\) generic coordinates. The
differential calculation gives \(\dim_L\Omega_{L/k}=d\). The \(dx_j\)
span this space: derivatives of their polynomial and rational
expressions follow from Leibniz and the quotient rule. Some \(d\) of
them therefore form a basis. Polynomial extension and localization, or
directly the universal property of derivations, give

\[
\Omega_{M/k}=(M\otimes_L\Omega_{L/k})
\oplus\bigoplus_{a,b}M\,dU_{ab}.
\]

Modulo the second summand, each \(dt_j\) is the corresponding generic
linear combination of the \(dx_h\). Their determinant in a differential
basis is a nonzero polynomial in the independent last-row coefficients:
specialize those rows to select the \(d\) basis generators to obtain
determinant one. Hence the \(dt_j\), together with all the \(dU_{ab}\),
form a basis of \(\Omega_{M/k}\). Normalization has already proved that
\(K=F(t)\) is a rational parameter field and \(M/K\) is finite. The
first fundamental sequence identifies \(\Omega_{M/K}\) with the quotient
by the span of these basis elements, so it is zero. Criterion (5a)
proves that \(M/K\) is separable. In characteristic zero the
finite-field calculation gives the same conclusion. The parameter field
\(F\) need not be perfect. \(\square\)

Put \(F_i=F(z_{i+1},\ldots,z_n)\), for \(0\le i\le n\), and let

\[
G_i=I_F F_i[z_1,\ldots,z_i]\cap F[z].
\]

These are the saturations by nonzero tail polynomials. A purely
transcendental field extension preserves primality and primaryness. For
the latter assertion, inject a primary quotient into its Artinian
localization at its radical. In the polynomial ring over this local
algebra, a polynomial with nonzero residue acts injectively: filter
coefficients by powers of the nilpotent maximal ideal; on the first
nonzero layer of another polynomial, multiplication is by a nonzero
polynomial over a field and cannot kill that vector-valued polynomial.
Polynomials with zero residue are nilpotent. Contraction proves
primaryness before this localization, and adjoining then localizing the
independent coefficient variables preserves it. Faithful flatness
preserves irredundancy. Thus the original primary decomposition and
associated primes extend faithfully to this generic coefficient field.
Lemma 5.1 and the primary localization rule show that \(G_i\) is the
intersection of precisely the primary components whose radicals have
dimension at least \(n-i\). This is a downward-closed set in prime
inclusion, so the contraction is canonical even when embedded primary
components are not unique. The ideals form
\(G_0\supset G_1\supset\cdots\supset G_n=I_F\). If \(I=0\), then all
\(G_i=0\); we treat that prime ideal separately. For nonzero \(I\),
\(G_0=F[z]\).

For \(1\le i\le n\), set

\[
H_i=(G_{i-1}/G_i)\otimes_{F[z_{i+1},\ldots,z_n]}F_i.
\tag{9}
\]

This is a finite-dimensional \(F_i\)-vector space, with commuting
multiplication operators from the \(z_j\). Indeed it embeds, by primary
decomposition, into the sum of the localized quotients for components of
dimension exactly \(n-i\); each is zero-dimensional over \(F_i\), hence
finite. Components of larger dimension impose \(G_{i-1}\) and give zero
in this injection; those of smaller dimension are killed by
localization. Every associated prime of dimension \(n-i\) actually
appears in \(H_i\): at that prime, irredundancy supplies an element of
the intersection of all other components which is absent from its
component; it is in \(G_{i-1}\) and survives in (9). In particular no
embedded dimension disappears from this filtration.

Adjoin further independent variables \(v_1,\ldots,v_n\), and on
\(H_i\otimes_{F_i}F_i(v)\) consider multiplication by
\(L_v=\sum_jv_jz_j\). Denote its monic minimal polynomial by
\(E_i(T,v)\) and its characteristic polynomial by \(R_i(T,v)\). For a
zero layer both are one. These are the elementary-divisor and norm
polynomials in this generic-fibre convention. Primitive numerator forms
can be obtained by clearing denominators; irreducibility below always
means over \(F_i(v)\), not an accidental factor introduced in a
denominator. This convention expresses the same generic
minimal-polynomial versus determinant mechanism as the Hentzelt--Noether
forms, while avoiding claims of equality between different coordinate
normalizations of the historical formulas.

\subsection{6. Prime and primary tests by the two
forms}\label{prime-and-primary-tests-by-the-two-forms}

\textbf{Theorem 6.1.} Let \(I\ne0\) be proper.

\begin{enumerate}
\def\labelenumi{\arabic{enumi}.}
\tightlist
\item
  The ideal is primary exactly when one layer is nonzero and its norm
  polynomial is a power of one irreducible polynomial. Equivalently its
  minimal polynomial is a power of one irreducible polynomial in that
  single layer.
\item
  A prime ideal has one layer and irreducible \(E_i\). Conversely, one
  layer with irreducible \(R_i\) implies that \(I\) is prime.
\item
  Over a perfect coefficient field, one layer with irreducible \(E_i\)
  is also equivalent to primality; for a prime ideal \(E_i=R_i\).
\item
  In characteristic \(p>0\), for a prime ideal over any field,
  \(R_i=E_i^{p^g}\) for some \(g\ge0\). Over an imperfect field an
  irreducible \(E_i\) can also occur for a proper primary ideal.
\end{enumerate}

The zero ideal is prime and is the separately excluded empty-filtration
case.

\textbf{Proof of the primary test.} Over an algebraic closure of
\(F_i\), the support points of a layer have distinct generic values
\(\sum v_j\xi_j\): equality as polynomials in independent \(v_j\) forces
equality of all coordinates. The determinant proof in Section 2, which
applies equally to finite modules supported at maximal ideals, makes its
characteristic polynomial the product of these linear factors with
positive multiplicities. Each residue field contributes one irreducible
factor over \(F_i(v)\), possibly with an inseparable power. Distinct
maximal ideals contribute different factors, because a common conjugate
generic value would identify their coordinate tuples and hence their
relation ideals. The minimal polynomial has the same distinct
irreducible factors as the characteristic polynomial: this follows on
each local module from scalar-plus-nilpotent multiplication, or from the
elementary-divisor decomposition of a linear operator.

By (9), the nonzero layers and their residue-field factors thus record
every associated prime of \(P/I\). One factor in one layer is equivalent
to one associated prime. A finite module with one associated prime is
primary: the zero divisors are the union of its associated primes, and
the radical of its annihilator is that single prime. Applied to \(P/I\)
this proves the first assertion, using
\href{../prerequisites/AG-CA/AG-CA-04.html}{Associated primes and
primary decomposition}, Theorems 1.2 and 4.2.

\textbf{Proof of the general prime assertions.} In the primary case the
only layer is the whole generic quotient
\(B=F_i[z_1,\ldots,z_i]/I_FF_i[z_1,\ldots,z_i]\), a finite local
algebra. Its residue field is denoted \(L\). If \(I\) is prime, \(B=L\)
is a field, so the minimal polynomial of \(L_v\) is irreducible. If the
characteristic polynomial is irreducible, its minimal polynomial equals
it. Then \(F_i(v)[L_v]\) is a field of degree \(\dim_{F_i}B\), and
therefore is all of \(B\otimes F_i(v)\). This algebra is a field, so
\(B\) is a field. A primary ideal contracts from its localization at
parameters outside its radical; consequently \(I_F\), and by faithful
field extension \(I\), is prime.

\textbf{Proof over a perfect field.} Lemma 5.1 makes \(L/F_i\)
separable. Its distinct embeddings give distinct values of the generic
\(L_v\); hence that value is a primitive element of \(L(v)/F_i(v)\). If
\(I\) is prime, this proves equality of minimal and characteristic
polynomials.

Now suppose the single-layer minimal polynomial \(h(T,v)\) is
irreducible. Reduction modulo the nilpotent maximal ideal of \(B(v)\)
gives the same minimal polynomial in its residue field, since that
polynomial is irreducible and annihilates the residue value. It is
separable by the preceding paragraph. Therefore \(h_T(L_v,v)\ne0\) in
the subfield \(F_i(v)[L_v]\subset B(v)\). Differentiate the identity
\(h(L_v,v)=0\) with respect to \(v_j\), keeping every element of \(B\)
fixed:

\[
h_T(L_v,v)z_j+h_{v_j}(L_v,v)=0.
\]

The first coefficient is invertible in that subfield. Hence every
\(z_j\) belongs to it. These elements generate \(B(v)\) as an algebra,
so \(B(v)\) is the field \(F_i(v)[L_v]\). The preceding contraction
argument proves primality.

\textbf{Proof of the inseparable assertion.} For a field \(L/F_i\), the
characteristic polynomial of a field element \(a\) is its minimal
polynomial to power \([L(v):F_i(v,a)]\): choose a basis over the
subfield and repeat its multiplication matrix. With \(a=L_v\), distinct
embeddings of \(L\) into an algebraic closure are distinguished by the
independent \(v_j\). Indeed all images of the \(z_j\) belong to an
algebraic closure of \(F_i\); an equality of their generic linear
combinations forces equality of each coordinate. The only embedding of
\(L(v)\) fixing \(F_i(v,L_v)\) is therefore the identity. The
finite-field calculation of Section 5 proves that this extension is
purely inseparable of degree \(p^g\), yielding the formula. \(\square\)

These are the modern forms of Noether's Theorems XII--XIV in Section 6.
``Proper primary'' there means primary but not prime. The perfect-field
hypothesis refers to the original coefficient field; the parameter field
itself may be imperfect. Generic separating parameters, not perfectness
of that rational function field, justify the proof.

For a concrete imperfect-field distinction, take \(k=\mathbb F_2(a,b)\),
with independent \(a,b\), and \(I=(x^2-a,y^2-b)\). The element \(a\) is
not a square, since its formal derivative with respect to \(a\) is one
and the derivative of every square is zero. After adjoining \(x\), the
field is \(\mathbb F_2(x,b)\), with \(a=x^2\); differentiating with
respect to \(b\) shows that \(b\) is still not a square. Thus the
quotient is a purely inseparable field of degree four. For
\(L_v=v_1x+v_2y\),

\[
E=T^2-v_1^2a-v_2^2b,\qquad R=E^2.
\]

The quadratic is irreducible since its constant term is not a square in
\(k(v)\): its derivative with respect to \(a\) is \(v_1^2\ne0\), whereas
every square has derivative zero. On the other hand, over
\(k=\mathbb F_2(a)\), the proper primary quotient
\(B=k[x,w]/(x^2-a,w^2)\) has the same pattern \(E=T^2-v_1^2a\),
\(R=E^2\). The nonzero nilpotent \(w\) proves that the ideal is not
prime, despite irreducible \(E\). These are the two mechanisms
distinguished by Noether's imperfect-field examples.

For a proper primary ideal the norm need not even be a \(p\)-power of
the irreducible minimal polynomial. Over \(\mathbb F_3(a)\), the
quotient by \((x^3-a,w^2)\) has basis \(x^j,x^jw\), \(0\le j<3\), and
therefore dimension six and length two over its degree-three residue
field. A generic value \(v_1x+v_2w\) has irreducible minimal polynomial
\(T^3-v_1^3a\): its cube is \(v_1^3a\), and differentiation with respect
to \(a\) proves that this coefficient is not a cube. The one-root
irreducible-polynomial calculation of Section 5 makes the degree three.
Its characteristic polynomial is the square of this polynomial, by the
composition-series argument. The exponent two is not a power of three.
The prime hypothesis in Theorem 6.1(4) is essential.

\subsection{7. What elimination
projects}\label{what-elimination-projects}

Let \(k\) be algebraically closed and write variables as \(x,y\). Set
\(J=I\cap k[x]\). The image of \(V(I)\) under projection to its
\(x\)-coordinates need not be closed, but its Zariski closure is exactly
\(V(J)\).

\textbf{Proof.} A polynomial in \(k[x]\) vanishes on that image exactly
when, viewed as a polynomial in all variables, it vanishes on \(V(I)\).
By the strong Nullstellensatz this ideal is
\(\sqrt I\cap k[x]=\sqrt J\); the equality follows directly by taking
powers. The closed set defined by the ideal of a set is its closure,
proving the assertion. \(\square\)

For \(I=(xy-1)\), the projection omits \(x=0\), yet \(J=0\) and the
closure is the whole line. Affine elimination equations determine a
closure, not necessarily the image itself.

For projective variables, closedness holds over every base, and
homogeneous elimination gives the actual image. Here is a full proof,
including the field and local algebra steps.

\textbf{Theorem 7.1 (projective projection and homogeneous
elimination).} For any scheme \(Y\), the projection
\(\mathbb P^r_Y\to Y\) sends closed subsets to closed subsets, and this
remains true after every change of base. On an affine base
\(Y=\operatorname{Spec}A\), let \(J\) be a homogeneous ideal of
\(S=A[X_0,\ldots,X_r]\), and let \(Z\) be its projective zero set. With
\(S_+=(X_0,\ldots,X_r)\), its image is

\[
\pi(Z)=V_A\bigl((J:S_+^\infty)\cap A\bigr),
\qquad
J:S_+^\infty=\bigcup_{N\ge0}(J:S_+^N).
\tag{10}
\]

Here \(J:S_+^N\) consists of the elements whose product with every
element of \(S_+^N\) is in \(J\). No Noetherian or finite-generation
hypothesis on \(A\) or \(J\) is needed.

\textbf{Proof of the projective field step.} Over a field \(K\), a
homogeneous ideal \(H\subset K[X_0,\ldots,X_r]\) has empty projective
zero set exactly when its quotient has zero degree-\(N\) piece for some
\(N\ge0\). Extend to an algebraic closure \(\bar K\). Empty projective
zero set means that all common affine zeros are contained in the origin.
Conversely a nonzero affine zero gives a projective point by scaling; if
the projective scheme has a point, one of its affine charts is a nonzero
finite-type \(K\)-algebra. A maximal ideal of that chart has finite
residue field by the weak Nullstellensatz, and an embedding of that
field in \(\bar K\) supplies a nonzero geometric zero. Thus this
geometric criterion detects all projective scheme points, not just
\(K\)-rational points.

The strong Nullstellensatz in \(\bar K[X]\) now puts each \(X_j\) in the
radical of \(H\bar K[X]\). Choose \(e_j\ge1\) with
\(X_j^{e_j}\in H\bar K[X]\). For \(N=1+\sum_{j=0}^r(e_j-1)\), every
monomial of degree \(N\) is divisible by at least one of those powers.
Hence the degree-\(N\) quotient is zero over \(\bar K\), and is already
zero over \(K\): tensoring a nonzero vector space with a field extension
cannot make it zero. A zero degree-\(N\) piece makes every higher piece
zero by multiplication by the variables, and forces all the \(X_j\) to
be nilpotent in the quotient. It therefore leaves no projective prime,
since a projective prime must avoid at least one \(X_j\). This proves
both directions. The argument also covers \(H=(1)\), where the
degree-zero piece is already zero.

\textbf{Proof near a point of the base.} Put \(M_N=(S/J)_N\). This is a
finite \(A\)-module, being a quotient of the free module on the finitely
many degree-\(N\) monomials. For a prime \(\mathfrak p\subset A\), its
fibre is

\[
M_N\otimes_A\kappa(\mathfrak p)
=\bigl(\kappa(\mathfrak p)[X]/J(\mathfrak p)\bigr)_N.
\]

Indeed taking this quotient and its fixed graded piece commutes with
tensor product by right exactness. Suppose the projective fibre is
empty. The field step gives an \(N\) for which this tensor product is
zero. Nakayama,
\href{../prerequisites/AG-CA/AG-CA-02.html}{Localization, local
properties and support}, Theorem 4.2, gives \((M_N)_{\mathfrak p}=0\).
Choose a finite generating list for \(M_N\). Each generator is killed by
some element outside \(\mathfrak p\), by the localization equality
criterion; their product \(f\notin\mathfrak p\) kills them all.
Consequently \((M_N)_f=0\). Multiplication by the variables gives
\((M_m)_f=0\) for every \(m\ge N\), so every projective fibre over the
open neighbourhood \(D(f)\) is empty. The complement of the image is
therefore open.

For the exact formula, a finite module \(M\) is zero after localization
at \(\mathfrak p\) exactly when \(\operatorname{Ann}_A(M)\) contains an
element outside \(\mathfrak p\): use the same finite generating list to
prove the nontrivial direction. The field step and Nakayama therefore
give

\[
\pi(Z)=\bigcap_{N\ge0}V_A(\operatorname{Ann}_A M_N)
=V_A\left(\bigcup_{N\ge0}\operatorname{Ann}_A M_N\right).
\]

The annihilators are increasing, because degree-\((N+1)\) monomials are
products of variables with degree-\(N\) monomials. Finally,
\(\operatorname{Ann}_A M_N=(J:S_+^N)\cap A\): both sides say that
multiplying every degree-\(N\) monomial by the coefficient gives an
element of \(J\), and those monomials generate \(S_+^N\) as an ideal.
This proves (10).

\textbf{Proof over general bases and under base change.} Every closed
subset \(W\) of \(\mathbb P^r_A\) is cut out by homogeneous equations.
To verify this topological assertion, take a point outside \(W\) in the
chart \(X_i\ne0\). An equation of \(W\) on that affine chart is nonzero
at the point. Homogenize it by clearing powers of \(X_i\), then multiply
by one further \(X_i\). The resulting homogeneous equation still does
not vanish at that point, vanishes on \(W\) within the chart, and
vanishes on all points outside the chart. The collection of these
equations cuts out exactly \(W\). The affine-base argument applies to
its homogeneous ideal. For arbitrary \(Y\), apply it on every affine
open: the intersection of the image with that open is the image of the
restricted closed subset, so it is closed there. Closedness is local on
this open cover. Finally, after any map \(Y'\to Y\), the pulled-back
projective space is \(\mathbb P^r_{Y'}\), to which the same proof
applies. This proves universal closedness. \(\square\)

In the classical case of an algebraically closed coefficient field,
formula (10) is the homogeneous elimination ideal on each affine chart
of the base. Saturation removes the irrelevant affine origin of the
projective coordinates. Thus projective elimination describes an image
that is already closed, while affine elimination may describe only its
closure.

Consider over \(\mathbb C\) the conics \(y^2-x=0\), \(y^2+x^2-2=0\).
Their resultant in \(y\) is

\[
(x^2+x-2)^2=(x-1)^2(x+2)^2.
\]

The points are \((1,\pm1)\) and \((-2,\pm i\sqrt2)\). The Jacobian
determinant is \(-2y(1+2x)\), nonzero at all four points, so each is
simple. The squares in the elimination resultant reflect two points
above each projected coordinate, rather than a double point at every
intersection. Here both polynomials are monic in the eliminated
variable, so no leading-coefficient exception interferes.

Higher-dimensional Chow forms organize incidence with complementary
linear spaces. They extend the theme of a form recording geometric
intersections; we make no additional Chow-form theorem a premise of the
proofs here.

\subsection{8. Exercises}\label{exercises}

\begin{enumerate}
\def\labelenumi{\arabic{enumi}.}
\tightlist
\item
  \textbf{Basic.} Compute \(\operatorname{Res}(x^2+bx+c,x-a)\) and
  explain its vanishing.
\item
  \textbf{Intermediate.} Compute the multiplication matrices for
  \((x^2-y,y^2-y)\), and check the characteristic polynomials in
  characteristic two as well as characteristic zero.
\item
  \textbf{Intermediate.} Compute the forms for \((x^2,y)\) and
  \((x^2-1,y-x)\). Identify what changes in characteristic two.
\item
  \textbf{Intermediate.} Eliminate \(y\) from \(x^2+y^2=1\), \(y=x^2-1\)
  over \(\mathbb C\), finding all points and local lengths.
\item
  \textbf{Advanced.} Prove the universal ideal-membership identity for
  the resultant using the adjugate, including degree bounds for its two
  coefficients. Explain why this proof survives specialization of
  leading coefficients.
\end{enumerate}

\subsection{9. Solutions}\label{solutions}

\textbf{1.} Formula (1) gives \(f(a)=a^2+ba+c\); the two sign changes
from \(\prod(\alpha_i-a)\) cancel. It vanishes exactly when \(a\) is a
root of the quadratic.

\textbf{2.} The basis and matrices in Section 4 follow from \(y=x^2\),
\(x^4=x^2\). Direct matrix multiplication gives \(M_y=M_x^2\).
Determinants give \(T^2(T^2-1)\) and \(T^2(T-1)^2\) over \(\mathbb Z\),
so reduction modulo any characteristic is valid. In characteristic two
both become \(T^2(T-1)^2\). The two length-two local factors yield these
multiplicities, as (4) requires.

\textbf{3.} The first quotient has basis \(1,x\), with \(M_x^2=0\),
\(M_y=0\), hence form \(u_0^2\). The second is \(k[x]/(x^2-1)\), with
\(y=x\), giving \(R=u_0^2-(u_1+u_2)^2\). In characteristic different
from two this factors as \((u_0+u_1+u_2)(u_0-u_1-u_2)\), representing
the simple points \((1,1),(-1,-1)\). In characteristic two it is the
square of the first linear form, representing one point of length two.

\textbf{4.} The monic linear relation substitutes directly into the
circle polynomial. The resultant is \(x^2+(x^2-1)^2-1=x^2(x^2-1)\). The
quotient is \(\mathbb C[x]/(x^2(x-1)(x+1))\), so the points are
\((0,-1),(1,0),(-1,0)\), with lengths \(2,1,1\). The repeated root at
zero is tangency, now a genuine local multiplicity rather than two
separate points with the same projection.

\textbf{5.} Form the universal matrix \(S\) over \(C\). The vector
\(w=\operatorname{adj}(S)e\), where \(e\) represents one, has entries in
\(C\). Its first \(n\) entries specify a polynomial of degree less than
\(n\), the remaining \(m\) a polynomial of degree less than \(m\). The
matrix identity \(Sw=(\det S)e\) is exactly
\(uf+vg=\operatorname{Res}(f,g)\). Every specialization is a ring
homomorphism preserving this identity. It uses no inverted coefficient;
only the interpretation of determinant vanishing as an actual common
root can fail after degrees drop.

\subsection{The historical papers}\label{the-historical-papers}

Read works 22 and 24 in the freely accessible
\href{https://zenodo.org/records/21923146/preview/00_EMMY_NOETHER_EN_COMPLETE_LINKED_READER.pdf?include_deleted=0}{complete
English corpus edition}, with its
\href{https://raw.githubusercontent.com/KokunoYumeto/emmy-noether-en/main/source/Noether_English_ED0014.tex}{editable
TeX}. Work 22, \emph{On the Theory of Polynomial Ideals and Resultants},
Sections 2--6, treats modules, norm and elementary-divisor forms. Work
24, \emph{Elimination Theory and General Ideal Theory}, Sections 5--6,
places the dimensional pieces in primary decomposition and states the
prime and proper-primary criteria as Theorems XII--XIV. The edition is a
machine-assisted working translation; its original and third-party
materials retain their own rights. The independent formulation above
explains our generic-fibre convention and proves its imperfect-field
qualifications; it does not identify every modern determinant with every
historical normalization.

\subsection{References}\label{references}

\begin{itemize}
\tightlist
\item
  K. Hentzelt, edited by Emmy Noether, \emph{On the Theory of Polynomial
  Ideals and Resultants}, original publication in Mathematische Annalen
  \textbf{88} (1923), 53--79; work 22, Sections 2--6, in the
  \href{https://zenodo.org/records/21923146/preview/00_EMMY_NOETHER_EN_COMPLETE_LINKED_READER.pdf?include_deleted=0}{free
  English corpus edition}.
\item
  Emmy Noether, \emph{Elimination Theory and General Ideal Theory},
  original publication in Mathematische Annalen \textbf{90} (1923),
  229--261; work 24, Sections 5--6 and Theorems XII--XIV, in the
  \href{https://zenodo.org/records/21923146/preview/00_EMMY_NOETHER_EN_COMPLETE_LINKED_READER.pdf?include_deleted=0}{same
  free edition}.
\item
  \href{../prerequisites/AG-CA/AG-CA-16.html}{Kähler differentials},
  Theorem 1.1, Theorems 2.2, 3.1 and 5.1, and Theorem 6.2, for the fully
  proved differential construction and rules. The additional finite
  separability criterion, perfect-field dimension calculation and
  generic separating parameters are proved in Section 5 here.
\item
  \href{../prerequisites/AG-CA/AG-CA-06.html}{The Nullstellensatz and
  Jacobson rings}, Theorems 1.3 and 2.1--2.2;
  \href{../prerequisites/AG-CA/AG-CA-02.html}{Localization, local
  properties and support}, Theorem 4.2. Together with the graded-piece
  argument in Theorem 7.1, these supply a complete internal proof of
  projective closedness.
\end{itemize}

\subsection{Editable sources}\label{editable-sources}

\href{ideal-theory-in-rings-source.zip}{Complete source archive} ·
\href{ideal-theory-in-rings.tex}{Complete course in LaTeX} ·
\href{NOE-IDL-05.tex}{This lesson in LaTeX}. The archive includes all
five lesson texts, the figures and their reproducible source.

\end{document}
